%%%%%%%%%%%%%%%%%%%%%%%%%%%%%%%%%%%%%%%%%%%%%%%%%%%%%%%%%%%%%%%%%%%%%%%%%%%%%%%%%%%%%%%%%%%%%%%%%%%%%%
%
%   Filename    : chapter_1.tex 
%
%   Description : This file will contain your Research Description.
%                 
%%%%%%%%%%%%%%%%%%%%%%%%%%%%%%%%%%%%%%%%%%%%%%%%%%%%%%%%%%%%%%%%%%%%%%%%%%%%%%%%%%%%%%%%%%%%%%%%%%%%%

\chapter{Introduction}
\label{sec:intro}    %--note: labels help you with hyperlink editing (using your IDE)

Dense small‑object counting and detection are essential tasks in public health, agriculture, and medical diagnostics. Mosquito‑borne diseases affect hundreds of millions of people each year, making vector surveillance a priority for disease prevention \cite{Javed2023}. Monitoring \textit{Aedes} mosquito populations relies on counting eggs collected from ovitraps to estimate breeding‑site activity and guide control measures \cite{NaranjoAlcazar2024}. In crop science, accurate counting of sorghum grains and seedlings supports early yield forecasting and speeds breeding programs, yet manual counting remains labor‑intensive and time‑consuming \cite{Chen2024,Santiago2024}. In digital pathology, automated detection of microscopic features in duodenal biopsies aids celiac disease diagnosis, where pathologists must discern subtle histologic patterns across many tissue patches \cite{Sali2019,Stoleru2022}. These applications demand high accuracy, but their reliance on parameter‑heavy architectures motivates a systematic investigation of model compression for efficient edge deployment.

\begin{sloppypar}
Recent work shows that domain‑specific detectors achieve strong performance, but all employ large‑scale models that require substantial memory and compute. EggCountAI \cite{Javed2023} attains 98.88\% egg‑counting accuracy with a Mask RCNN containing millions of parameters. Sorghum‑grain counting \cite{Bakshi2025,Freeman2023,Santiago2024} uses heavyweight YOLO variants trained on GPU workstations. State‑of‑the‑art YOLO‑based approaches reach expert‑level performance for multi‑class cell detection and counting in celiac histology patches \cite{Tyagi2023}. Although compression techniques \cite{Endrawati2021,Wei2018} such as pruning have accelerated generic detection networks, pruning alone often degrades small‑object detection accuracy. To address this, recent studies combine pruning with knowledge distillation (KD) and quantisation. KD restores accuracy lost during pruning by transferring feature knowledge from a large teacher model to a compact student, improving small‑object detection in models like ScaleKD \cite{zhu2023scalekd} and PKD‑YOLOv8 \cite{Yu2025_PKD}. The approach also proves effective in medical imaging, where frameworks such as FKDNuSeg employ KD to create highly compact models for nuclei segmentation \cite{Zhao2026_FKDNuSeg}. Low‑bit quantisation further speeds inference on edge devices by converting floating‑point numbers to 8‑bit or 4‑bit integers. Methods like Reg‑PTQ \cite{Ding2024_RegPTQ} and Q‑DETR \cite{Xu2023_QDETR} specifically reduce quantisation errors in bounding‑box coordinates and transformer queries, ensuring dense, small objects remain accurately detected \cite{Ding2024_RegPTQ, Xu2023_QDETR}.
\end{sloppypar}

The present work proposes a systematic, multi‑technique compression methodology that reduces parameter counts and computational latency while preserving small‑object detection performance. Unlike prior efforts that focus on a single model, dataset, or isolated compression strategy, this study evaluates the integration of structured pruning, knowledge distillation, and quantisation across multiple detection architectures, including YOLO‑based models and RT‑DETR. Experiments span several extreme‑density datasets covering entomology, agriculture, and pathology. Results reveal consistent patterns in module sensitivity, identify optimal compression configurations that deliver substantial hardware acceleration with minimal performance loss, and offer practical guidelines for deploying parameter‑efficient small‑object detectors across diverse domains.

\section{Background of the Study}
\label{sec:overview}

Deep learning has dramatically improved computer vision, enabling highly accurate automated object detection and counting. State‑of‑the‑art models, such as single‑stage convolutional networks and transformer‑based detectors, recognise complex visual patterns with precision. However, these models are typically large and require significant memory and floating‑point operations, limiting their deployment to high‑performance GPU workstations. Model compression is therefore essential for deploying such systems on resource‑constrained edge devices \cite{He2025,Wei2018}. Structured pruning, knowledge distillation, and network quantisation aim to systematically reduce model parameters and latency while preserving acceptable performance on target hardware.

Within this context, dense small‑object detection and counting are critical in several fields. Public‑health vector surveillance depends on counting microscopic mosquito eggs collected from field traps \cite{Javed2023,NaranjoAlcazar2024}. Agricultural monitoring requires counting thousands of densely packed sorghum grains or seedlings to estimate crop yield \cite{Chen2024,Santiago2024}. In digital pathology, automated cell detection on histology patches supports clinical diagnoses \cite{Sali2019,Stoleru2022,Tyagi2023}. These small targets occupy few pixels and are often densely packed, placing high demands on a detector's multi‑scale feature extraction and spatial representation capabilities.

Most existing methods rely on heavyweight, uncompressed detectors to ensure high accuracy \cite{Bakshi2025,Zhao2023RTDETR}. Although structured pruning can reduce model size, applying it alone often degrades the feature‑aggregation necks and attention layers vital for small‑object localisation \cite{Endrawati2021}. Recent research therefore combines pruning with other compression techniques. Knowledge distillation (KD) restores lost accuracy by transferring feature representations from a large teacher network to a smaller student, proving effective for pest detection \cite{Yu2025_PKD,Zhu2023_ScaleKD} and nuclei segmentation \cite{Zhao2026_FKDNuSeg}. Low‑bit network quantisation maps high‑precision parameters to integer formats, accelerating on‑device execution while preserving regression performance \cite{Ding2024_RegPTQ,Xu2023_QDETR}.

A major challenge remains in applying these techniques to dense counting tasks. Prior work typically evaluates structured pruning, knowledge distillation, and quantisation in isolation, or limits assessment to generic datasets. No systematic, cross‑domain study has examined how a combined, multi‑step compression pipeline influences the distinct structural modules of YOLO and RT‑DETR architectures under extreme target density. The absence of a unified, architecture‑aware optimisation framework hampers the ability to balance model size, edge inference latency, and counting accuracy across diverse application domains.

\section{Research Objectives}

The general objective of this research is to design and evaluate a multi‑step model compression framework for the resource‑constrained edge deployment of parameter‑efficient small‑object detectors.

The specific objectives are:
\begin{enumerate}
    \item Analyse baseline performance for various detection architectures across diverse high‑density domains.
    \item Develop a multi‑step compression pipeline to systematically reduce model complexity.

    \item Evaluate the trade‑offs between model compression ratios, inference latency, and overall detection reliability.
\end{enumerate}

\section{Scope and Limitations of the Research}
\label{sec:scopelimitations}

The first objective establishes a robust baseline for performance, computational efficiency, and parameter profiles across high‑density domains. The scope is limited to three specialised small‑object detection datasets: the \textit{Aedes aegypti} egg dataset \cite{vicente2024aedesaegyptieggcounting}, the CoNSeP dataset \cite{Tyagi2023}, and the Sorghum‑Grain‑Count (SGC) dataset \cite{Bakshi2025}. These datasets represent high object density in biological, medical, and agricultural domains. All inputs are resized to $640\times640$ pixels, preserving the train/validation/test splits defined in the original papers.

The second objective develops and evaluates a compression pipeline employing structured pruning, knowledge distillation, and quantisation. Structured pruning removes redundant filters from single or multiple architectural layers, while knowledge distillation transfers feature knowledge from larger models to compact students. Quantisation reduces the precision of weights. The scope confines the study to these methods to assess their collective impact on preserving spatial features and counting accuracy without extensive architectural redesign.

The third objective assesses the feasibility of deploying optimised models on edge hardware. Testing is limited to on‑device inference on edge hardware platforms, without any off‑device processing or cloud uploads. This validates the practical utility of the compressed models in isolated, resource‑constrained environments where real‑time, on‑device analysis is required.

The fourth objective measures trade‑offs between compression ratios, inference latency, and detection reliability. Evaluation focuses on quantitative metrics, specifically the relationship between parameter reduction and mean average precision (mAP). This establishes a measurable Pareto front, providing a framework for selecting optimal compression configurations based on deployment constraints.

\section{Significance of the Research}
\label{sec:significance}

This research advances computer‑vision by addressing the gap between the computational demands of high‑precision object detection and the hardware limitations of edge‑computing environments. By creating a multi‑step compression framework that integrates structured pruning, knowledge distillation, and quantisation, the study supplies empirical evidence on how these techniques interact to preserve feature‑extraction fidelity under substantial parameter reduction. The contribution offers a systematic methodology for maintaining counting accuracy in high‑density domains such as biological, medical, and agricultural settings, where large‑scale architectures are often computationally prohibitive. The framework provides the community with a scalable template for deploying advanced detection models on resource‑constrained hardware while retaining performance reliability.

Furthermore, the work demonstrates the viability of on‑device intelligence. Enabling high‑performance inference on consumer‑grade mobile devices and single‑board computers eliminates the latency and privacy concerns associated with cloud‑dependent architectures. Stakeholders in agriculture and public health can thus obtain real‑time autonomous diagnostics in field conditions lacking reliable network infrastructure. By facilitating edge‑deployed counting tools, this research lays a technical foundation for more accessible and privacy‑preserving data analysis, reducing economic and operational barriers to implementing computer‑aided monitoring in critical industries.
